\section{서론}

웹 그래픽 사용자 인터페이스(GUI)는 텍스트, 레이아웃, 현재 상태,
가능한 행동을 함께 해석해야 하는 텍스트 밀집 시각 환경이다. 최근에는
DOM과 규칙 기반 도구에 의존하는 에이전트에서 벗어나, 스크린샷을
직접 보고 브라우저 행동을 예측하는 시각언어 에이전트가 활발히
연구되고 있다 \citep{qin2025uitars,gupta2026molmoweb}. 한편
VisualWebBench는 현재 멀티모달 모델이 웹페이지 OCR, 세부 이해,
그라운딩에서 여전히 큰 성능 격차를 보인다는 점을 드러냈다
\citep{liu2024visualwebbench}. 대규모 웹 UI 명령어 데이터가 웹
과제뿐 아니라 문서와 차트 등 비웹 텍스트 밀집 과제에도 전이된다는
결과 \citep{liu2025multiui}는 웹페이지 supervision이 범용적인 시각
표현을 학습하는 데 유용할 가능성을 보여준다.

그러나 기존 웹 GUI 모델은 데이터 출처, 과제 혼합 비율, 행동 공간,
학습 단계가 서로 달라 개별 설계 선택의 효과를 분리하기 어렵다.
특히 새로운 기반 모델에서 웹 GUI 학습을 시작할 때
\emph{명령 수행이나 행동 궤적을 배우기 전에 웹페이지로부터 무엇을
학습시켜야 하는가}라는 기본 질문에 명확한 답이 없다. LLaVA 계열과
같은 범용 멀티모달 모델은 이미지--텍스트 정렬 또는 시각 명령어
튜닝을 먼저 수행한다 \citep{liu2023llava,liu2024llavanext}. 웹 GUI
모델도 지각 및 그라운딩 데이터를 사용하지만, 화면 설명의 적절한
상세도는 분명하지 않다.

짧은 설명은 페이지의 목적, 핵심 상태, 대표 상호작용에 supervision을
집중한다. 반대로 긴 설명은 작은 텍스트와 상태, 표의 각 항목,
행동 가능한 요소까지 보존하여 더 풍부한 웹 표현을 만들 수 있다.
그러나 장문은 중요도가 낮은 요소의 나열에 모델 용량을 소비하거나,
후속 명령어 및 궤적 학습에서 그 효과가 사라질 수도 있다. 그림
\ref{fig:contrast}는 본 연구의 두 supervision 조건과 공통으로
통제되는 요소를 요약한다.

\begin{figure}[t]
\centering
\begin{tikzpicture}[
  font=\footnotesize,
  box/.style={draw, rounded corners=2pt, align=center,
              minimum height=8mm, inner sep=2.5pt},
  arrow/.style={-{Latex[length=1.8mm]}, thick},
  note/.style={font=\scriptsize, align=center}
]
\node[box, fill=stagegray, text width=.78\columnwidth] (screen)
  {\textbf{동일한 웹 화면과 Stage 1 evidence}\\
   content · state · relation · affordance · coordinate};
\node[box, fill=shortblue, below=7mm of screen,
      xshift=-.27\columnwidth, text width=.25\columnwidth] (concise)
  {\textbf{Concise}\\salient evidence};
\node[box, fill=stagegreen, below=7mm of screen,
      text width=.25\columnwidth] (balanced)
  {\textbf{Balanced}\\section-level\\evidence};
\node[box, fill=longorange, below=7mm of screen,
      xshift=.27\columnwidth, text width=.25\columnwidth] (comprehensive)
  {\textbf{Compre-}\\\textbf{hensive}\\near-complete\\evidence};
\node[box, fill=stagegray, below=10mm of balanced,
      text width=.78\columnwidth] (eval)
  {\textbf{공통 grounding 및 후속 학습}\\
   coverage 분석 · targeted diagnostic\\공개 benchmark · agent transfer};
\draw[arrow] (screen.south) -- (concise.north);
\draw[arrow] (screen.south) -- (balanced.north);
\draw[arrow] (screen.south) -- (comprehensive.north);
\draw[arrow] (concise.south) -- (eval.north west);
\draw[arrow] (balanced.south) -- (eval.north);
\draw[arrow] (comprehensive.south) -- (eval.north east);
\end{tikzpicture}
\caption{\textbf{웹페이지 verbalization의 정보 포괄성 비교.}
동일한 화면과 Stage 1 evidence를 서로 다른 information budget으로
언어화한다. 세 recipe는 평균 coverage가 증가하지만 sample별
evidence 집합의 엄격한 포함관계를 가정하지 않는다. Grounding
데이터와 후속 학습은 동일하게 유지한다.}
\label{fig:contrast}
\end{figure}


본 연구는 \textbf{간결한 설명(Concise)}과 \textbf{포괄적
설명(Comprehensive)}을 비교한다. 두 조건은 스크린샷, 1단계 구조화
사실, 좌표 규약, 외부 그라운딩 데이터, 기반 모델, 최적화 설정,
후속 데이터를 공유한다. 바뀌는 것은 설명의 길이와 이에 수반되는
정보 선택뿐이다. 효과는 사전학습(PT) 직후, 공통 명령어 튜닝(IT)
직후, 행동 궤적 지도학습(Traj) 직후의 세 지점에서 측정한다.

본 연구의 기여는 다음과 같다.
\begin{itemize}
    \item 그라운딩 양과 후속 데이터를 고정한 상태에서 웹 GUI
    사전학습의 서술 상세도를 비교하는 통제 실험을 설계한다.
    \item 아홉 개 체크포인트를 통해 PT의 효과가 어느 단계에서
    나타나고 후속 학습 이후까지 유지되는지 추적한다.
    \item 접근성 트리의 bounding box 좌표를 teacher가 추정하지 않고
    인용하도록 하며, 프로그램 검증으로 좌표 발명을 제거하는 2단계
    합성 파이프라인을 제시한다.
\end{itemize}
현재 주요 평가가 진행 중이므로 어느 조건이 우월하다고 미리
주장하지 않고, 기여를 연구 질문과 방법론 중심으로 기술한다.
